\section{江淮、两浙与四川：一条防线背后的几个南宋}

\subsection{从扬州的仓廪到临安的案牍}

绍兴初年的朝廷在扬州、越州、杭州之间移动时，最难搬运的并非仪仗，而是能让一纸命令变成粮食、舟船和人手的关系。北方失守后沿江各地接到的，不是一份已经配齐资源的重建方案；州县要安顿过江的官员、宗室、军人和流民，又要保留本地仓储、渡船和耕作。高宗在建炎元年即位，\eventref{SS-1127-EVENT-111}{南宋重建}由此开始，可是“朝廷已立”并不等于税粮已经可以按时抵达。扬州的驻跸、杭州的兵变和金军渡江，依次显示了行在对地方军队、河道和城门的依赖。

《建炎以来系年要录》把诏奏、行幸和任免排入连续日期，读来仿佛中枢始终有一张清楚的行动表；它真正能够可靠显示的是中枢如何理解自己的困局。每一次改驻、调将或催运，都要经过已有的水路、仓场、地方吏员和能够垫付的人家。地方是否照数交纳、船户怎样避开战线、逃难者在何处落籍，不能由这部编年书一笔断定。《宋史》本纪保留的政局节点与《要录》相互照应，却同样把镜头更多留在行在。把两种材料同渡口、城址和后来的地方文献并读，至少能看见国家不是从杭州向四方均匀铺开，而是在若干可供给的节点之间重新接线。

临安后来成为行在，\eventref{SS-1138-EVENT-118}{它所代表的不是一座旧都城的复制}。钱塘江口、运河和太湖流域使官署、市场、漕粮与海运较易相接，也使这座城市承受了更多驻军、官员和消费人口。城市的繁盛不应倒过来解释为南宋天然富强：临安能吸纳人和货物，恰因为东南乡村、盐场、织作地、平码船和转运节点被纳入一个持续征调的网络。市面上可以见到的货物与官府账簿上的收入不是同一件事；前者说明交换仍在发生，后者说明某个机关试图把一部分交换分类、征收和运走。

制度得失也在这里显出两面。将行在固定在临安，解决了朝廷反复迁徙时文书与决策难以积累的问题；代价则是国防的重心更依赖江淮和长江水道。江南的富庶被用来购买时间，时间又要用军粮、军器、俸给和赈济来维持。后来每一次北伐或守江的争论，表面是战略选择，实际都要回答同一个较难的问题：哪一段河道、哪一批转运人、哪一种税收和信用可以承担新的远征。

\subsection{淮河不是一根画在地图上的线}

绍兴和议以后，淮河、大散关一带常被概括为“宋金边界”。对驻防者和居民而言，这条边界首先是一串城、寨、渡口、堤岸和道路。两淮的城镇需要在农时、运粮和警报之间分配劳力；渡口既能运兵，也能运盐、米和逃亡者；同一条水道在不同季节会给守方或攻方不同的便利。和议降低了大规模会战的频率，却没有让边区退出财政。驻军的粮饷、马匹、器械、修城和信使，仍要从州县和转运体系取得。

因此，绍兴和议\eventref{SS-1141-EVENT-120}{不能被读成战争的句号}。它给朝廷以整顿军政的间隔，也把较长期的军费、使节和边市安排制度化。宋方叙事常把屈辱、恢复或和战的道德语言放在前面；金方材料则有自己的战功与外交措辞。两者相合之处可以确认谈判和边防确曾持续，彼此相异之处提醒读者：盟书的称谓不能直接告诉我们一座边城实际由谁征税，一条路上又有多少家庭迁离。兵额和岁币的数字尤其不能脱离报送机关、时间和用途，变成可以横量“国力”的单一刻度。

一一六一年金军南下，采石一带的战事使长江防线再次成为现实的运输问题。江面宽阔并不自动保护南宋；船只、风向、江岸停泊处、守军能否接到补给，都决定渡江企图的成本。采石之后的局势变化，和高宗禅位、孝宗即位接续在一起。新朝廷面对的是一支需要饷给的军队、一个已经习惯和议的边境以及希望恢复北方的政治期待。孝宗北伐在符离受挫，随后订立隆兴和议，不能化为“主战必败”或“主和必懦”的格言；远征所需的道路与协同没有自动随命令出现，失败又使前线和江南承担下一轮整军、转运和议和的成本。

两淮社会不应只在战报里出现。围城、迁军和征粮会改变市镇开闭、田地能否播种和家庭如何避险，但能留下的文字极不均匀。墓志和地方碑刻容易保存官员、士绅与寺院的行动，官府条文则保存应当如何安置、赈济或禁约；它们不能替无名佃户和船户说出完整经历。正因为如此，叙事只能把可见的军政节点同水道、城镇和征发压力连接，而不虚构一座城中所有人的选择。

\subsection{四川不是东南财政的附图}

南宋的西部防线有另一种地理。四川的山地、峡江和盆地使道路、栈道、江运与关隘彼此咬合；它能给防御者以纵深，也会让调兵、送粮和传递消息更加困难。朝廷从东南调度资源时，不能把四川看作一块可以用同一张账簿管理的地方。当地军政、盐井、山城、市场和移民社会有各自的节奏，战时的命令抵达与实际执行之间尤其容易出现时间差。

十二世纪末到十三世纪初，金宋关系和边防形势屡变，四川既要面对北方压力，也要处理内地供给。南宋的会要和食货材料保存了机构如何设置、何种赋役或军需被讨论的线索；这些条目首先证明政府提出过什么办法，不足以证明办法已经覆盖川峡每一处山谷。地方文书、碑刻、城址与水利遗存能让某些地点更具体，但它们又各有保存范围。将一座山城的坚守写成全蜀社会的共同意志，将一件契约写成普遍的土地关系，都会把材料的局部性误作总体。

蒙古向南推进时，四川的空间问题更突出。钓鱼城相关战事之所以持续进入宋、元和蒙古叙事，不只是因为一位统帅的生死传说，而是因为嘉陵江流域的城防、山道、粮运和区域控制彼此牵连。传世叙事对蒙哥死因和战斗细节并不一致；能够较稳妥地说的是一二五九年前后的战事改变了蒙古对宋战争的节奏，而不能把每一个后出的说法都当作城上所见。考古的城址分期可以约束防御空间的想象，却不能独自复原每日攻守或全部居民的生活。

四川的制度得失在于：以山江形势和地方军政维持了一段难以迅速突破的纵深，代价是补给链长、区域差异大，中央越在战争中需要资源，越要经由本地中介。这样一种防御并不能简单归为强或弱；它把攻击者的成本推高，也使南宋很难把东南的财政能力直接换成西部每一处防线的稳定。

\subsection{江南的钱粮如何成为边防}

“江南财赋”是南宋叙事里最容易说得过于轻松的词。它不是一只可以随时打开的国库，而是一连串收取、折纳、储存、运输、支付和核销的动作。田租、盐、茶、商税、榷货和临时征发进入不同机关，可能以米、绢、钱或票据的形式流转；到达边区之前还要经过仓场、船队、脚夫、文书和地方关系。任何一个环节受洪水、盗贼、战事、价格或官吏更替影响，前线得到的便不再是诏书上那个数目。

会子的发展把这种困难的一部分转化为信用和票据问题。纸币可以减轻某些携运负担，使中央较快地支付和调拨；它也要求发行、回收、折换和市场信任能够配合。会要中反复出现的限额、回收和管理措施，恰好说明制度需要不断调整，不能被读作一种已经无摩擦运行的货币机器。\eventref{SS-1150-EVENT-122}{会子管理}所能显示的是朝廷试图处理财政工具；它不能使我们从发行额直接算出临安、两淮或四川一家人的购买力，更不能把物价变化归于单一政策。

财政压力还通过地方社会分配。能够预缴、承运或与官府有关系的家族，和靠短工、佃作、渔盐或小额贸易维生的人，面对的风险并不相同；但现存材料更常记录有名有产者的契约和功德。南宋的契约、墓志和地方材料为研究者保留了少数地区的具体关系，价值正在于它们让抽象的“赋役”落到某处田宅、某个债务或一次运输上。它们不能被扩大为全境统计。书写这种不对称，并非逃避解释，而是防止东南的材料密度把边区和无名者再一次挤出历史。

\subsection{地区差异不是国家失败的旁注}

到十三世纪后期，襄樊和长江中游的战争把几个地区重新连到一起。围城需要封锁汉水与陆路，守方需要让江南的资源沿水道北上；四川的战局和荆襄道路又影响援军能否相互接应。襄阳、樊城失守，\eventref{SS-1273-EVENT-145}{不是一座孤城的结局}，而是多年供给、舰船、城防和区域调度逐渐失去平衡的结果。它同时提醒我们，东南市场的活力不能自动跨越被截断的通道，城防的坚固也不能代替一条持续的运输线。

临安在一二七六年降元，\eventref{SS-1276-EVENT-148}{标志的是朝廷中枢的转换}；沿江、闽广和海上发生的事情仍有不同步的时间。对有些地方而言，接收意味着新的官府和税役；对另一些人而言，首先是逃难、改行或与新旧军队谈判。宋末与元初的文本容易把人分为忠、降、叛等整齐类别，地方材料却常只能让我们看见一次捐施、一份迁葬记录或一段断裂的经营。读者应把这些片段放回具体地点，不把后来形成的道德图谱当作所有人的内心。

南宋在区域差异中维持了一个多世纪。它的成就不是把所有地方塑造成同一制度的复制品，而是在江海、边河、山城和市场之间保持足以征收、运输、裁判和防守的连接。它的副作用也正在这里：连接越依赖富庶的东南与少数水路节点，断裂的代价就越集中。下一节转向港口和市镇，所要追问的不是“商业是否繁荣”这一抽象问题，而是船、货、税、寺院与法庭怎样在战争国家里共同塑造可用的资源。
